\documentclass[conference,10pt]{IEEEtran}
\usepackage{cite}
\usepackage{amsmath,amssymb}
\usepackage{graphicx}
\usepackage{url}
\usepackage{booktabs}
\usepackage{listings}
\usepackage{multirow}
\usepackage{caption}
\usepackage{subcaption}
\hyphenation{op-tical net-works semi-conduc-tor}

\title{Design and Implementation of Efficient Job Scheduling and Data Structures for Process Management}

\author{
  \IEEEauthorblockN{<STUDENT1_NAME>\IEEEauthorrefmark{1}, <STUDENT2_NAME>\IEEEauthorrefmark{1}, <STUDENT3_NAME>\IEEEauthorrefmark{1}, <STUDENT4_NAME>\IEEEauthorrefmark{1}, <STUDENT5_NAME>\IEEEauthorrefmark{1}, Prof. <PROF_NAME>\IEEEauthorrefmark{2}} 
  \IEEEauthorblockA{\IEEEauthorrefmark{1}Department of Computer Science, <INSTITUTE_NAME>\\
    Emails: <student1@email>, <student2@email>, <student3@email>, <student4@email>, <student5@email>}
  \IEEEauthorblockA{\IEEEauthorrefmark{2}Department of Computer Science, <INSTITUTE_NAME>\\
    Email: <prof.email@institute.edu>}
  \thanks{Manuscript received Month DD, 20XX; revised Month DD, 20XX.}
}

\begin{document}

\maketitle

\begin{abstract}
This paper presents the design, implementation, and evaluation of a compact job-management system that demonstrates efficient scheduling and supporting data structures for process control. The system integrates a priority queue, min--max heap, trie-based job correction/indexing, and process control block (PCB) management to support prioritized job scheduling, fast lookup, and real-time corrections. We describe the algorithms used, analyze their time and space complexities, and evaluate the system on representative workloads. Results show that the combined approach yields competitive performance for insertion, deletion, and lookup operations while providing predictable scheduling behavior. % Replace with measured results and key metrics.
\end{abstract}

\begin{IEEEkeywords}
Job scheduling, priority queue, min--max heap, trie, process control block, algorithm analysis, data structures.
\end{IEEEkeywords}

\section{Introduction}
The efficient management of jobs and processes is central to operating systems and many real-time applications. This work implements and evaluates a compact scheduler leveraging several classical and hybrid data structures to balance fast lookup, correction, and prioritized dispatch. Contributions include: (1) implementation of an array-based priority queue API, (2) a min--max heap for double-ended priority access, (3) a trie-based JobCorrector for fast key correction and suggestions, and (4) an integrated PCB structure tying scheduling metadata to queue operations. The remainder of the paper is organized as follows.

\section{Background and Motivation}
Scheduling policies such as FCFS, SJF, Priority, and Round-Robin offer different trade-offs in response time, throughput, and fairness. For workloads requiring prioritized service and fast textual lookup/correction of job identifiers (e.g., command names, job keys), combining priority-based queues with trie-based indexing offers practical benefits: O(log n) priority operations and O(L) key operations (L = key length).

\section{Literature Review}
This section surveys foundational work and related approaches relevant to our design choices, grouped by topic.

\subsection{Priority Queues and Heaps}
Priority queues are a core abstract data type used in scheduling and many algorithms. Binary heaps provide a simple array-backed implementation with \(O(\log n)\) insertion and extraction and \(O(1)\) peek, and are discussed in-depth in algorithm textbooks \cite{cormen2009introduction}. For workloads that require efficient access to both minima and maxima, min--max heaps were introduced by Atkinson et al. \cite{atkinson1986minmax}; they maintain alternating min and max levels in a binary heap layout and support both delete-min and delete-max in \(O(\log n)\).

Fibonacci heaps provide amortized \(O(1)\) insert and decrease-key operations and \(O(\log n)\) extract-min, useful in specialized contexts such as Dijkstra's algorithm \cite{fredman1987fibonacci}. However, their constant factors and implementation complexity often make binary heaps preferable in practice for scheduling tasks with frequent inserts and extracts.

\subsection{Double-Ended and Hybrid Priority Structures}
Double-ended priority queues and specialized structures (e.g., pairing heaps, min--max heaps) allow efficient retrieval of both extremes. Prior work compares designs for different workload characteristics and operation mixes; Atkinson et al.'s min--max heap remains one of the simplest practical solutions \cite{atkinson1986minmax}.

\subsection{Trie Structures and String Correction}
Tries (prefix trees) are a classical structure for associative arrays keyed by strings; Fredkin's early work \cite{fredkin1960trie} formalized the idea. For approximate matching and spell-correction, techniques such as Levenshtein automata, BK-trees, and edit-distance-aware traversals are commonly used \cite{navarro2001approximate, gusfield1997algorithms}. For small alphabets and high locality in keys, compact trie variants (compressed tries, radix trees) reduce space while preserving \(O(L)\) query time \cite{fredkin1960trie, gusfield1997algorithms}.

\subsection{Process Control Blocks and Scheduler Design}
Operating-system-level scheduler design and PCB representations are covered in standard texts \cite{silberschatz2018operating, tanenbaum2015modern}. The PCB typically stores identifiers, state, scheduling priority, arrival time, and resource accounting. Our implementation uses a lightweight PCB struct that links job metadata with queue positions and supports easy state transitions (ready, running, blocked, terminated).

\subsection{Recent Practical and Hybrid Approaches}
Recent engineering-focused work often favors simpler, well-understood data structures (binary heap, binary search tree variants, tries) with careful engineering to minimize cache misses and memory overhead \cite{weiss2014data}. Hybrid approaches combine structures to meet multi-objective constraints (e.g., fast approximate lookup + prioritized dispatch), similar in spirit to our design. When extreme decrease-key performance is required, Fibonacci heaps or pairing heaps have been explored, but with caution due to complexity and overhead \cite{fredman1987fibonacci}.

\subsection{Summary of Gaps Addressed}
While prior work provides mature structures for priority management and string indexing separately, there is less emphasis on small educational schedulers that combine trie-based correction/indexing with priority/min--max queue dispatch. This project demonstrates an integrated approach, evaluates operation-level costs, and documents trade-offs relevant for moderate-scale scheduling tasks.

\section{System Design and Methodology}
% (Same high-level content as template: architecture, data structures, algorithms.)

\subsection{Data Structures}
\textbf{Priority Queue:} array-based binary heap with index mapping for decrease-key (if needed). Complexity: insert \(O(\log n)\), extract \(O(\log n)\), peek \(O(1)\).

\textbf{Min--Max Heap:} binary layout with alternating min and max levels, offering both delete-min and delete-max in \(O(\log n)\). Useful when both extremes are relevant for scheduling policies.

\textbf{Trie (JobCorrector):} standard pointer-based trie storing job keys and job-ID lists at terminal nodes; lookup and insertion are \(O(L)\) where \(L\) is the key length. For approximate correction, we implement single-edit heuristics and prefix suggestions; more advanced approaches (Levenshtein automaton) are noted for future work \cite{navarro2001approximate}.

\textbf{PCB:} a small struct with job ID, priority, arrival time, runtime estimate, and current state. PCBs are referenced from queue entries to update job metadata without costly searches.

\section{Algorithms and Complexity}
This section presents clear pseudocode for the core operations and a formal complexity analysis.

\subsection{Priority Queue (binary heap)}
Pseudocode (array-backed max-heap):
\begin{lstlisting}[language=C]
// insert(key)
heap[++size] = key;
idx = size;
while (idx > 1 && heap[idx] > heap[idx/2]) {
  swap(heap[idx], heap[idx/2]);
  idx = idx/2;
}

// extract_max()
ret = heap[1];
heap[1] = heap[size--];
idx = 1;
while (2*idx <= size) {
  child = 2*idx;
  if (child+1 <= size && heap[child+1] > heap[child]) child++;
  if (heap[idx] >= heap[child]) break;
  swap(heap[idx], heap[child]);
  idx = child;
}
return ret;
\end{lstlisting}

Complexity: insertion and extract operations cost $O(\log n)$ time due to heapify up/down; peek is $O(1)$. Space is $O(n)$ for the array.

\subsection{Min--Max Heap}
The min--max heap alternates levels: even levels maintain min-heap order and odd levels maintain max-heap order. Core operations (insert, delete-min, delete-max) perform at most $O(\log n)$ comparisons and swaps. Implementation follows Atkinson et al.'s level-based bubble operations \cite{atkinson1986minmax}.

\subsection{Trie (JobCorrector)}
Pseudocode for insertion and exact/approximate lookup (single-edit heuristic):
\begin{lstlisting}[language=C]
// insert(root, key)
node = root;
for (c in key) {
  if (!node->child[c]) node->child[c] = new TrieNode();
  node = node->child[c];
}
node->isTerminal = true; node->jobList.add(jobId);

// lookup with single-edit suggestions (DFS up to edit budget)
search(node, i, editsLeft)
  if (i == len(key) && node->isTerminal) collect(node)
  // match next char
  for (ch in alphabet) {
    if (node->child[ch]) {
      search(node->child[ch], i + (ch==key[i]?1:0), editsLeft - (ch==key[i]?0:1))
    }
  }
\end{lstlisting}

Exact lookup is $O(L)$ where $L$ is the key length. Approximate (single-edit) search scales as $O(A\cdot L)$ where $A$ is alphabet size times branching caused by allowed edits; in practice we bound edits to 1 which keeps the search practical.

\subsection{Scheduling Loop}
High-level scheduler pseudocode:
\begin{lstlisting}[language=C]
while (!queue.empty()) {
  job = priority_queue.extract_max();
  dispatch(job);
  update_pcb(job);
  if (job.needs_correction) corrected = trie_suggest(job.key);
}
\end{lstlisting}

Per-job dispatch overhead is dominated by the priority queue extract ($O(\log n)$) and any optional trie-based correction ($O(L)$ or bounded approximate cost).

\subsection{Aggregate Complexity}
For a workload of $m$ job arrivals of average key length $\bar L$:
- Building/inserting all jobs: $m\cdot O(\log n)$
- Serving all jobs: $m\cdot O(\log n + \bar L)$ if corrections are requested for each job.

Space complexity: $O(n + S_{trie})$ where $n$ is number of queued jobs and $S_{trie}$ is trie memory (dependent on total characters stored).

\section{Implementation}
The implementation is in C. Main source files map to repository components such as `PriorityQueue.c`/`PriorityQueue.h`, `MinMax.c`/`MinMaxHeap.h`, `JobTrie.c`/`JobTrie.h`, `JobCorrector.c`/`JobCorrector.h`, `PCB.c`/`PCB.h`, and `main.c` driver harness. Build flags used for experiments: `-std=c11 -O2 -Wall`.

\section{Experimental Setup and Results}
We ran microbenchmarks and end-to-end traces to evaluate operation latency and scheduler throughput. Replace or extend these placeholders with concrete measurements from your test harness.

\subsection{Workloads}
- Synthetic: uniformly random job keys, priority drawn from a Zipf/uniform distribution.
- Realistic: replay traces or curated sequences with bursts and varying key lengths.

\subsection{Metrics}
- Throughput (jobs / second)
- Average wait time (ms)
- Operation latencies: insert, extract, lookup (median and 95th percentile)
- Memory usage (peak resident set size)

\subsection{Representative Results (example placeholders)}
\begin{table}[ht]
\caption{Example operation latencies (replace with measured values)}
\centering
\begin{tabular}{lrr}
\toprule
Operation & Median (\micro s) & 95th (\micro s) \\
\midrule
Insert & 120 & 340 \\
Extract & 140 & 410 \\
Trie lookup (exact) & 35 & 85 \\
Trie lookup (single-edit) & 420 & 920 \\
\bottomrule
\end{tabular}
\end{table}

Interpretation: the heap-based priority queue dominates per-dispatch overhead with $O(\log n)$ behavior; trie exact lookups are inexpensive but approximate searches have higher variability. Use these placeholders to insert your measured results and replace the example table.

\section{Discussion}
This section highlights trade-offs, limitations, and optimization opportunities.

\textbf{Trade-offs.} The design favors predictability and simplicity: binary and min--max heaps give good practical $O(\log n)$ performance with low implementation overhead. Tries give $O(L)$ exact lookup and useful prefix/correction functionality but incur extra memory for nodes. When keys share prefixes, memory usage is amortized; otherwise node counts grow with character volume.

\textbf{Limitations.} The current implementation is single-threaded and does not account for concurrent producers/consumers. Decrease-key is supported via index mappings but not optimized for extreme workloads; in those cases specialized heaps (Fibonacci/pairing) may help but increase complexity.

\textbf{Optimizations.}
- Use a compressed trie (radix/Patricia trie) to reduce node overhead.
- Use memory pools/arena allocators for trie nodes to reduce allocation overhead and improve cache locality.
- For heavy decrease-key workloads consider pairing heaps or Fibonacci heaps only if decrease-key dominates.
- Implement lock-free or fine-grained locking for concurrent schedulers.

\textbf{Correctness and Robustness.} Unit tests should verify heap invariants after randomized operations and validate trie insertion/lookup including approximate suggestions. Instrumentation (counters, assertions) helps detect corruption early.

\section{Conclusion}
We presented a compact job-management system that integrates array-based priority queues, a min--max heap, trie-based correction/indexing, and a lightweight PCB representation. The implementation emphasizes clear, practical algorithms with predictable $O(\log n)$ scheduling behavior and $O(L)$ key operations. The integrated design is suitable for educational use and moderate-scale scheduling tasks where fast lookups and correction suggestions are required. Future empirical evaluation (scale-up experiments and real-trace replay) will quantify trade-offs more precisely.

\section{Future Work}
Planned extensions include:
- Concurrency: support multiple producer/consumer threads with lock-free or fine-grained locking for queues and tries.
- Persistence: back the job queue and PCB store with on-disk structures for restart/resume.
- Advanced correction: integrate Levenshtein automata or ML-based suggestion models for higher-quality corrections at scale.
- Benchmarks: run larger-scale experiments (millions of jobs) and compare against alternative heap implementations (pairing, Fibonacci) and database-backed queues.
- Visualization: add tracing and visualization utilities to profile queue behavior and key-level correction patterns.

\section*{Acknowledgments}
% (Optional acknowledgments.)

\bibliographystyle{IEEEtran}
\begin{thebibliography}{99}
\bibitem{cormen2009introduction}
T. H. Cormen, C. E. Leiserson, R. L. Rivest, and C. Stein, \textit{Introduction to Algorithms}, 3rd ed., MIT Press, 2009.

\bibitem{atkinson1986minmax}
M. D. Atkinson, J.-R. Sack, N. Santoro, and I. Stolfi, "Min--max heaps and generalized priority queues," \textit{Communications of the ACM}, vol. 29, no. 10, pp. 996--1000, 1986.

\bibitem{fredkin1960trie}
E. Fredkin, "Trie memory," \textit{Communications of the ACM}, vol. 3, no. 9, pp. 490--499, 1960.

\bibitem{silberschatz2018operating}
A. Silberschatz, P. B. Galvin, and G. Gagne, \textit{Operating System Concepts}, 10th ed., Wiley, 2018.

\bibitem{tanenbaum2015modern}
A. S. Tanenbaum and H. Bos, \textit{Modern Operating Systems}, 4th ed., Pearson, 2015.

\bibitem{fredman1987fibonacci}
M. L. Fredman and R. E. Tarjan, "Fibonacci heaps and their uses in improved network optimization algorithms," \textit{Journal of the ACM}, vol. 34, no. 3, pp. 596--615, 1987.

\bibitem{gusfield1997algorithms}
D. Gusfield, \textit{Algorithms on Strings, Trees and Sequences: Computer Science and Computational Biology}, Cambridge University Press, 1997.

\bibitem{navarro2001approximate}
G. Navarro, "A guided tour to approximate string matching," \textit{ACM Computing Surveys}, vol. 33, no. 1, pp. 31--88, 2001.

\bibitem{weiss2014data}
M. A. Weiss, \textit{Data Structures and Algorithm Analysis in C}, 3rd ed., Pearson, 2014.

\bibitem{tycoon1999pairing}
R. E. Tarjan and M. L. Fredman, "Pairing heaps and potential function analysis," (classic pairing heap references and analyses). % Replace with full citation if used.

% Add any additional papers, articles, or online resources consulted during implementation.
\end{thebibliography}

% Appendix: Example compile command
% gcc -std=c11 -O2 -Wall main.c PriorityQueue.c MinMax.c JobTrie.c JobCorrector.c PCB.c -o scheduler

\end{document}
